On considère une seconde solution aqueuse $\mathrm{(S^{\prime})}$ d'acide
propanoïque $\ce{C2H5COOH}$ de concentration molaire
$C_{A}^{\prime}=\qty{0,10}{\M}$. La valeur du taux d'avancement final de la
réaction de l'acide propanoïque avec l'eau est $\tau^{\prime}=\num{1,16e-3}$.
\begin{enumerate}
    \item En comparant $\tau^{\prime}$ avec $\tau$ le taux d'avancement final de
          la réaction d'acide méthanoïque avec l'eau, indiquer lequel des deux
          acides est le plus dissocié en solution.
    \item Comparer les constantes d'acidité
          $\mathrm{K_{A}}(\ce{HCOOH_{(aq)}}/\ce{HCOO^-_{(aq)}})$ et
          $\mathrm{K_{A}}(\ce{C2H5COOH_{(aq)}}/\ce{C2H5COO^-_{(aq)}})$.
\end{enumerate}

